%% kapa: exam=1 round=34 year=2023 period=1 form=A subject=real_estate subject_ko=부동산학원론 seq=17 number=97
%% source: Q-Net 감정평가사 기출문제 (cst00309) artlSeq=5211789 file=raw/1008/5211789/2023년도 제34회 감정평가사 1차시험 1교시.pdf
%% conversion: kordoc 4.24.51 -> tools/gen1.py (auto); review=proofread-sonnet; flags=none
\begin{question}{97}
부동산투자에 관한 설명으로 옳은 것은?
\begin{choices}
  \item 부동산투자는 부동산이 갖고 있는 고유한 특성이 있지만 환금성, 안전성 측면에서 주식 투자와 다르지 않다.
  \item 부동산은 실물자산이기 때문에 인플레이션 방어 능력이 우수하여 디플레이션과 같은 경기침체기에 좋은 투자대상이다.
  \item 부동산은 다른 투자상품에 비하여 거래비용의 부담이 크지만 부동산시장은 정보의 대칭성으로 인한 효율적 시장이다.
  \item 부동산투자는 부동산의 사회적·경제적·행정적 위치의 가변성 등으로 인해 부동산시장의 변화를 면밀히 살펴야 한다.
  \item 투자의 금융성이란 투자자가 투자자산을 필요한 시기에 손실없이 현금으로 전환할 수 있는 안전성의 정도를 말한다.
\end{choices}
\end{question}
\answer{4}
